% main.tex
% gs -dNOPAUSE -dBATCH -sDEVICE=png16m -sOutputFile=out%03d.png -r1000x1000 main.pdf
% ffmpeg -framerate 30 -i out%03d.png out.mp4
% https://tex.stackexchange.com/a/748667
\unprotect
    %%% Defined similarly to \obeyMPboxdepth, \ignoreMPboxdepth,
    %%% \obeyMPboxorigin, and \normalMPboxdepth
    \permanent\protected\def\strutMPboxdepth{%
        \let\meta_relocate_box\meta_strut_box_depth%
    }

    \protected\def\meta_strut_box_depth{%
        \setbox\b_meta_graphic\vpack{%
            \box\b_meta_graphic%
            %%% Add some depth to the box without affecting its size
            \nohrule height-\openstrutdepth width 0pt depth \openstrutdepth\relax%
        }
    }

    %%% The pagecolumns output routine alredy adds a \vfil elsewhere, so the one
    %%% here isn't needed and disrupts the depth added by \strutMPboxdepth.
    \let\page_otr_fill_page_unchecked\relax
\protect

%%% Comment the line below to confirm that an overfull vbox is generated;
%%% uncomment it to confirm that the overfull vbox message goes away.
\strutMPboxdepth
\setupbodyfont[pagella]


\definepapersize
  [youtubeHD]
  [width=1920bp,height=1080bp]
\setuppapersize[youtubeHD][youtubeHD]
\setuplayout[
    location=middle,
    width=middle,
    height=middle,
    topspace=0pt,
    backspace=0pt,
    header=0pt,
    footer=0pt
]
\defineframed[myframed][
    width=800pt,
    align=normal,
    offset=1cm,
    frame=off
]
\setupbodyfont[modern,24pt]
\ctxloadluafile{parametric.lua}
\processMPfigurefile{register_mp_cmd.mp}

\startMPinclusions
def render_theorem_box =
numeric w, h;
w := CurrentWidth;
h := CurrentHeight;

path p;
p := (-w/2, -h/2) -- (w/2, -h/2) -- (w/2, h/2) -- (-w/2, h/2) -- cycle;
setbounds currentpicture to p;

picture theorem;
theorem := textext("\getbuffer[BufA]");

pair center_anchor, topleft_anchor;
center_anchor := center theorem;
topleft_anchor := ulcorner theorem;

enddef;


def jasper_thin =
    withpen pencircle scaled .1cm
enddef ;

def jasper_thick =
    withpen pencircle scaled .2cm
enddef ;

def jasper_lightgray =
    withcolor (0.6)
enddef ;

def jasper_lightgraypenthin =
    jasper_lightgray withpen pencircle scaled .1cm
enddef ;

def jasper_lightgraypenthick =
    jasper_lightgray withpen pencircle scaled .2cm
enddef ;


def jasper_darkgray =
    withcolor (0.4)
enddef ;

def jasper_darkgraypenthin =
    jasper_darkgray withpen pencircle scaled .15cm
enddef ;

def jasper_darkgraypenthick =
    jasper_darkgray withpen pencircle scaled .3cm
enddef ;
\stopMPinclusions


\starttext

    \startbuffer[BufA]
        \startframed[myframed]
            {\bf Theorem 3.2.1} There exists an angle congruent to
            its supplement.\par 
            \color[white]{{\it Proof.} Let \im{AB} be a straight line and \im{M}
            be a point not lying on \im{AB}. Let us subject the 
            plane to a type (iii) isometry: reflection in \im{AB}.
            That means the isometry will carry every point on one 
            side of \im{AB} into a point lying on the other side 
            of this line, and every point of \im{AB} will remain 
            stationary.\par 
            As a result of this transformation, point \im{M} will 
            be carried into some point \im{N} lying on the other 
            side of \im{AB}.\par
            Let us join \im{M} with \im{N} by means of a straight 
            line. This line will intersect \im{AB} at 
            some point, which we shall denote \im{C}.\par 
            Since \im{CA} does not move as a result of the isometry
            performed, \im{\angle MCA} is transformed into 
            \im{\angle NCA}, which means that these two angles 
            are congruent: \im{\angle MCA = \angle NCA}. At the 
            same time, since the two sides (\im{MC} and \im{CN})
            of these angles extend each other into a straight line
            (\im{MN}), these angles are supplementary. Therefore
            \im{\angle MCA} and \im{\angle NCA} are congruent 
            supplementary angles, which means such angles exist,
            {\bf Q.E.D.}}
        \stopframed
    \stopbuffer
    % initial theorem
    \dorecurse{7*30}{%
        \startMPpage
            render_theorem_box ; 
            numeric t;
            t := 0;
            pair anchor;
            anchor := (1 - t) * center_anchor + t * topleft_anchor + t * (-2cm, -9cm);
            label(theorem, (0,3cm) - anchor);
        \stopMPpage%
    }
    % Translate the theorem to the top right corner.
    \dorecurse{2*30+1}{%
        \startMPpage
            render_theorem_box ;
            numeric t;
            t := (#1-1)/60;
            pair anchor;
            anchor := (1 - t) * center_anchor + t * topleft_anchor + t * (-2cm, -9cm);
            label(theorem, (0,3cm) - anchor);
        \stopMPpage%
    }%
    \startMPinclusions
        def my_text =
            render_theorem_box ;
            numeric t;
            t := 1; 
            pair anchor;
            anchor := (1 - t) * center_anchor + t * topleft_anchor + t * (-2cm, -9cm);
            label(theorem, (0,3cm) - anchor);
        enddef ;
    \stopMPinclusions%
    % Pause once the theorem is in place.
    \dorecurse{4*30}{%
        \startMPpage
            my_text ;
        \stopMPpage%
    }
    \startbuffer[BufA]
        \startframed[myframed]
            {\bf Theorem 3.2.1} There exists an angle congruent to
            its supplement.\par 
            {\it Proof.} \color[white]{Let \im{AB} be a straight line and \im{M}
            be a point not lying on \im{AB}. Let us subject the 
            plane to a type (iii) isometry: reflection in \im{AB}.
            That means the isometry will carry every point on one 
            side of \im{AB} into a point lying on the other side 
            of this line, and every point of \im{AB} will remain 
            stationary.\par 
            As a result of this transformation, point \im{M} will 
            be carried into some point \im{N} lying on the other 
            side of \im{AB}.\par
            Let us join \im{M} with \im{N} by means of a straight 
            line. This line will intersect \im{AB} at 
            some point, which we shall denote \im{C}.\par 
            Since \im{CA} does not move as a result of the isometry
            performed, \im{\angle MCA} is transformed into 
            \im{\angle NCA}, which means that these two angles 
            are congruent: \im{\angle MCA = \angle NCA}. At the 
            same time, since the two sides (\im{MC} and \im{CN})
            of these angles extend each other into a straight line
            (\im{MN}), these angles are supplementary. Therefore
            \im{\angle MCA} and \im{\angle NCA} are congruent 
            supplementary angles, which means such angles exist,
            {\bf Q.E.D.}}
        \stopframed
    \stopbuffer
    % {\it Proof.}
    \dorecurse{3*30}{%
        \startMPpage
            my_text ;
        \stopMPpage%
    }
    \startbuffer[BufA]
        \startframed[myframed]
            {\bf Theorem 3.2.1} There exists an angle congruent to
            its supplement.\par 
            {\it Proof.} Let \im{AB} be a straight line and \im{M}
            be a point not lying on \im{AB}. \color[white]{Let us subject the 
            plane to a type (iii) isometry: reflection in \im{AB}.
            That means the isometry will carry every point on one 
            side of \im{AB} into a point lying on the other side 
            of this line, and every point of \im{AB} will remain 
            stationary.\par 
            As a result of this transformation, point \im{M} will 
            be carried into some point \im{N} lying on the other 
            side of \im{AB}.\par
            Let us join \im{M} with \im{N} by means of a straight 
            line. This line will intersect \im{AB} at 
            some point, which we shall denote \im{C}.\par 
            Since \im{CA} does not move as a result of the isometry
            performed, \im{\angle MCA} is transformed into 
            \im{\angle NCA}, which means that these two angles 
            are congruent: \im{\angle MCA = \angle NCA}. At the 
            same time, since the two sides (\im{MC} and \im{CN})
            of these angles extend each other into a straight line
            (\im{MN}), these angles are supplementary. Therefore
            \im{\angle MCA} and \im{\angle NCA} are congruent 
            supplementary angles, which means such angles exist,
            {\bf Q.E.D.}}
        \stopframed
    \stopbuffer
    % Let \im{AB} be a straight line and \im{M}
    % be a point not lying on \im{AB}.
    \dorecurse{10*30}{%
        \startMPpage
            my_text ;
        \stopMPpage%
    }
    \startMPinclusions
        def A_B =
            my_text ;
            jasper_point[
                fx="cx"
                ,fy="cy+7"
            ] ;
            jasper_point[
                fx="cx"
                ,fy="cy-7"
            ] ;
            jasper_append_label[
                px="cx+cos(5*pi/4)"
                ,py="cy+7+sin(5*pi/4)"
                ,text="\im{A}"
            ] ;
            jasper_append_label[
                px="cx+cos(5*pi/4)"
                ,py="cy-7+sin(5*pi/4)"
                ,text="\im{B}"
            ] ;
        enddef ;
    \stopMPinclusions
    % generate points A and B
    \dorecurse{3*30}{%
        \startMPpage
            A_B ;
            jasper_render_segments ;
        \stopMPpage%
    }
    % draw line AB
    \dorecurse{2*30}{%
        \startMPpage
            A_B ;
            jasper_curve[
                ustart="-9"
                ,ustop="9"
                ,usamples="200"
                ,fx="cx"
                ,fy="cy+u*#1/60"
                ,drawcommands="jasper_thin"
            ] ;
            jasper_render_segments ;
        \stopMPpage%
    }%
    \startMPinclusions
        def lineAB =
            A_B ;
            jasper_curve[
                ustart="-9"
                ,ustop="9"
                ,usamples="20"
                ,fx="cx"
                ,fy="cy+u"
                ,drawcommands="jasper_thin"
            ] ;
        enddef ;
    \stopMPinclusions
    % pause after drawing line AB
    \dorecurse{2*30}{%
        \startMPpage
            lineAB ;
            jasper_render_segments ;
        \stopMPpage%
    }%
    \startMPinclusions
        def point_M =
            lineAB ;
            jasper_point[
                fx="cx-10"
                ,fy="cy-1"
            ] ;
            jasper_append_label[
                px="cx-10+cos(5*pi/4)"
                ,py="cy-1+sin(5*pi/4)"
                ,text="\im{M}"
            ] ;
        enddef ;
    \stopMPinclusions%
    % draw point M, exterior to AB
    \dorecurse{5*30}{%
        \startMPpage
            point_M ;
            jasper_render_segments ;
        \stopMPpage%
    }%
    \startbuffer[BufA]
        \startframed[myframed]
            {\bf Theorem 3.2.1} There exists an angle congruent to
            its supplement.\par 
            {\it Proof.} Let \im{AB} be a straight line and \im{M}
            be a point not lying on \im{AB}. Let us subject the 
            plane to a type (iii) isometry: reflection in \im{AB}.
            \color[white]{That means the isometry will carry every point on one 
            side of \im{AB} into a point lying on the other side 
            of this line, and every point of \im{AB} will remain 
            stationary.\par 
            As a result of this transformation, point \im{M} will 
            be carried into some point \im{N} lying on the other 
            side of \im{AB}.\par
            Let us join \im{M} with \im{N} by means of a straight 
            line. This line will intersect \im{AB} at 
            some point, which we shall denote \im{C}.\par 
            Since \im{CA} does not move as a result of the isometry
            performed, \im{\angle MCA} is transformed into 
            \im{\angle NCA}, which means that these two angles 
            are congruent: \im{\angle MCA = \angle NCA}. At the 
            same time, since the two sides (\im{MC} and \im{CN})
            of these angles extend each other into a straight line
            (\im{MN}), these angles are supplementary. Therefore
            \im{\angle MCA} and \im{\angle NCA} are congruent 
            supplementary angles, which means such angles exist,
            {\bf Q.E.D.}}
        \stopframed
    \stopbuffer
    % Let us subject the 
    % plane to a type (iii) isometry: reflection in \im{AB}.
    \dorecurse{10*30}{%
        \startMPpage
            point_M ;
            jasper_render_segments ;
        \stopMPpage%
    }%
    \startbuffer[BufA]
        \startframed[myframed]
            {\bf Theorem 3.2.1} There exists an angle congruent to
            its supplement.\par 
            {\it Proof.} Let \im{AB} be a straight line and \im{M}
            be a point not lying on \im{AB}. Let us subject the 
            plane to a type (iii) isometry: reflection in \im{AB}.
            That means the isometry will carry every point on one 
            side of \im{AB} into a point lying on the other side 
            of this line, and every point of \im{AB} will remain 
            stationary.\par 
            \color[white]{As a result of this transformation, point \im{M} will 
            be carried into some point \im{N} lying on the other 
            side of \im{AB}.\par
            Let us join \im{M} with \im{N} by means of a straight 
            line. This line will intersect \im{AB} at 
            some point, which we shall denote \im{C}.\par 
            Since \im{CA} does not move as a result of the isometry
            performed, \im{\angle MCA} is transformed into 
            \im{\angle NCA}, which means that these two angles 
            are congruent: \im{\angle MCA = \angle NCA}. At the 
            same time, since the two sides (\im{MC} and \im{CN})
            of these angles extend each other into a straight line
            (\im{MN}), these angles are supplementary. Therefore
            \im{\angle MCA} and \im{\angle NCA} are congruent 
            supplementary angles, which means such angles exist,
            {\bf Q.E.D.}}
        \stopframed
    \stopbuffer
    % That means the isometry will carry every point on one 
    % side of \im{AB} into a point lying on the other side 
    % of this line, and every point of \im{AB} will remain 
    % stationary.
    \dorecurse{15*30}{%
        \startMPpage
            point_M ;
            jasper_render_segments ;
        \stopMPpage%
    }%
    \startbuffer[BufA]
        \startframed[myframed]
            {\bf Theorem 3.2.1} There exists an angle congruent to
            its supplement.\par 
            {\it Proof.} Let \im{AB} be a straight line and \im{M}
            be a point not lying on \im{AB}. Let us subject the 
            plane to a type (iii) isometry: reflection in \im{AB}.
            That means the isometry will carry every point on one 
            side of \im{AB} into a point lying on the other side 
            of this line, and every point of \im{AB} will remain 
            stationary.\par 
            As a result of this transformation, point \im{M} will 
            be carried into some point \im{N} lying on the other 
            side of \im{AB}.\par
            \color[white]{Let us join \im{M} with \im{N} by means of a straight 
            line. This line will intersect \im{AB} at 
            some point, which we shall denote \im{C}.\par 
            Since \im{CA} does not move as a result of the isometry
            performed, \im{\angle MCA} is transformed into 
            \im{\angle NCA}, which means that these two angles 
            are congruent: \im{\angle MCA = \angle NCA}. At the 
            same time, since the two sides (\im{MC} and \im{CN})
            of these angles extend each other into a straight line
            (\im{MN}), these angles are supplementary. Therefore
            \im{\angle MCA} and \im{\angle NCA} are congruent 
            supplementary angles, which means such angles exist,
            {\bf Q.E.D.}}
        \stopframed
    \stopbuffer
    % As a result of this transformation, point \im{M} will 
    % be carried into some point \im{N} lying on the other 
    % side of \im{AB}.
    \dorecurse{13*30}{%
        \startMPpage
            point_M ;
            jasper_render_segments ;
        \stopMPpage%
    }%
    % move N into position
    \dorecurse{3*30}{%
        \startMPpage
            point_M ;
            jasper_point[
                fx="cx-10+20*#1/90"
                ,fy="cy-1"
            ] ;
            jasper_append_label[
                px="cx-10+20*#1/90+cos(5*pi/4)"
                ,py="cy-1+sin(5*pi/4)"
                ,text="\im{N}"
            ] ;
            jasper_render_segments ;
        \stopMPpage%
    }%
    \startMPinclusions
        def point_N =
            point_M ;
            jasper_point[
                fx="cx-10+20"
                ,fy="cy-1"
            ] ;
            jasper_append_label[
                px="cx-10+20+cos(5*pi/4)"
                ,py="cy-1+sin(5*pi/4)"
                ,text="\im{N}"
            ] ;
        enddef ;
    \stopMPinclusions
    % point N
    \dorecurse{5*30}{%
        \startMPpage
            point_N ;
            jasper_render_segments ;
        \stopMPpage%
    }%
    \startbuffer[BufA]
        \startframed[myframed]
            {\bf Theorem 3.2.1} There exists an angle congruent to
            its supplement.\par 
            {\it Proof.} Let \im{AB} be a straight line and \im{M}
            be a point not lying on \im{AB}. Let us subject the 
            plane to a type (iii) isometry: reflection in \im{AB}.
            That means the isometry will carry every point on one 
            side of \im{AB} into a point lying on the other side 
            of this line, and every point of \im{AB} will remain 
            stationary.\par 
            As a result of this transformation, point \im{M} will 
            be carried into some point \im{N} lying on the other 
            side of \im{AB}.\par
            Let us join \im{M} with \im{N} by means of a straight 
            line. \color[white]{This line will intersect \im{AB} at 
            some point, which we shall denote \im{C}.\par 
            Since \im{CA} does not move as a result of the isometry
            performed, \im{\angle MCA} is transformed into 
            \im{\angle NCA}, which means that these two angles 
            are congruent: \im{\angle MCA = \angle NCA}. At the 
            same time, since the two sides (\im{MC} and \im{CN})
            of these angles extend each other into a straight line
            (\im{MN}), these angles are supplementary. Therefore
            \im{\angle MCA} and \im{\angle NCA} are congruent 
            supplementary angles, which means such angles exist,
            {\bf Q.E.D.}}
        \stopframed
    \stopbuffer
    % Let us join \im{M} with \im{N} by means of a straight 
    % line.
    \dorecurse{7*30}{%
        \startMPpage
            point_N ;
            jasper_render_segments ;
        \stopMPpage%
    }%
    % draw line MN
    \dorecurse{3*30}{%
        \startMPpage
            point_N ;
            jasper_curve[
                ustart="0"
                ,ustop="20"
                ,usamples="200"
                ,fx="cx-10+u*#1/90"
                ,fy="cy-1"
                ,drawcommands="jasper_thin"
            ] ;
            jasper_render_segments ;
        \stopMPpage%
    }%
    \startMPinclusions
        def lineMN =
            point_N ;
            jasper_curve[
                ustart="0"
                ,ustop="20"
                ,usamples="200"
                ,fx="cx-10+u"
                ,fy="cy-1"
                ,drawcommands="jasper_thin"
            ] ;
        enddef ;
    \stopMPinclusions%
    % pause after drawing MN
    \dorecurse{5*30}{%
        \startMPpage
            lineMN ;
            jasper_render_segments ;
        \stopMPpage%
    }%
    \startbuffer[BufA]
        \startframed[myframed]
            {\bf Theorem 3.2.1} There exists an angle congruent to
            its supplement.\par 
            {\it Proof.} Let \im{AB} be a straight line and \im{M}
            be a point not lying on \im{AB}. Let us subject the 
            plane to a type (iii) isometry: reflection in \im{AB}.
            That means the isometry will carry every point on one 
            side of \im{AB} into a point lying on the other side 
            of this line, and every point of \im{AB} will remain 
            stationary.\par 
            As a result of this transformation, point \im{M} will 
            be carried into some point \im{N} lying on the other 
            side of \im{AB}.\par
            Let us join \im{M} with \im{N} by means of a straight 
            line. This line will intersect \im{AB} at 
            some point, which we shall denote \im{C}.\par 
            \color[white]{Since \im{CA} does not move as a result of the isometry
            performed, \im{\angle MCA} is transformed into 
            \im{\angle NCA}, which means that these two angles 
            are congruent: \im{\angle MCA = \angle NCA}. At the 
            same time, since the two sides (\im{MC} and \im{CN})
            of these angles extend each other into a straight line
            (\im{MN}), these angles are supplementary. Therefore
            \im{\angle MCA} and \im{\angle NCA} are congruent 
            supplementary angles, which means such angles exist,
            {\bf Q.E.D.}}
        \stopframed
    \stopbuffer
    % This line will intersect \im{AB} at 
    % some point, which we shall denote \im{C}.
    \dorecurse{10*30}{%
        \startMPpage
            lineMN ;
            jasper_render_segments ;
        \stopMPpage%
    }%
    \startMPinclusions
        def pointC =
            lineMN ;
            jasper_point[
                fx="cx"
                ,fy="cy-1"
            ] ;
            jasper_append_label[
                px="cx+cos(5*pi/4)"
                ,py="cy-1+sin(5*pi/4)"
                ,text="\im{C}"
            ] ;
        enddef ;
    \stopMPinclusions%
    % point C
    \dorecurse{5*30}{%
        \startMPpage
            pointC ;
            jasper_render_segments ;
        \stopMPpage%
    }%
    \startbuffer[BufA]
        \startframed[myframed]
            {\bf Theorem 3.2.1} There exists an angle congruent to
            its supplement.\par 
            {\it Proof.} Let \im{AB} be a straight line and \im{M}
            be a point not lying on \im{AB}. Let us subject the 
            plane to a type (iii) isometry: reflection in \im{AB}.
            That means the isometry will carry every point on one 
            side of \im{AB} into a point lying on the other side 
            of this line, and every point of \im{AB} will remain 
            stationary.\par 
            As a result of this transformation, point \im{M} will 
            be carried into some point \im{N} lying on the other 
            side of \im{AB}.\par
            Let us join \im{M} with \im{N} by means of a straight 
            line. This line will intersect \im{AB} at 
            some point, which we shall denote \im{C}.\par 
            Since \im{CA} does not move as a result of the isometry
            performed, \im{\angle MCA} is transformed into 
            \im{\angle NCA}, \color[white]{which means that these two angles 
            are congruent: \im{\angle MCA = \angle NCA}. At the 
            same time, since the two sides (\im{MC} and \im{CN})
            of these angles extend each other into a straight line
            (\im{MN}), these angles are supplementary. Therefore
            \im{\angle MCA} and \im{\angle NCA} are congruent 
            supplementary angles, which means such angles exist,
            {\bf Q.E.D.}}
        \stopframed
    \stopbuffer
    % Since \im{CA} does not move as a result of the isometry
    % performed, \im{\angle MCA} is transformed into 
    % \im{\angle NCA},
    \dorecurse{12*30}{%
        \startMPpage
            pointC ;
            jasper_render_segments ;
        \stopMPpage%
    }%
    % build \im{\angle MCA}
    \dorecurse{4*30}{%
        \startMPpage
            pointC ;
            jasper_curve[
                ustart="0"
                ,ustop="pi/2"
                ,fx="cx+2*cos(pi-u*#1/120)"
                ,fy="cy-1+2*sin(pi-u*#1/120)"
                ,fz="1"
                ,drawcommands="jasper_lightgraypenthin"
            ] ;
            jasper_curve[
                ustart="0.8"
                ,ustop="1.2"
                ,fx="cx+u*2*cos(pi-(pi/4)*#1/120)"
                ,fy="cy-1+u*2*sin(pi-(pi/4)*#1/120)"
                ,fz="1"
                ,drawcommands="jasper_lightgraypenthin"
            ] ;
            jasper_render_segments ;
        \stopMPpage%
    }%
    % pause after building angle MCA
    \dorecurse{4*30}{%
        \startMPpage
            pointC ;
            jasper_curve[
                ustart="0"
                ,ustop="pi/2"
                ,fx="cx+2*cos(pi-u)"
                ,fy="cy-1+2*sin(pi-u)"
                ,fz="1"
                ,drawcommands="jasper_lightgraypenthin"
            ] ;
            jasper_curve[
                ustart="0.8"
                ,ustop="1.2"
                ,fx="cx+u*2*cos(pi-(pi/4))"
                ,fy="cy-1+u*2*sin(pi-(pi/4))"
                ,fz="1"
                ,drawcommands="jasper_lightgraypenthin"
            ] ;
            jasper_render_segments ;
        \stopMPpage%
    }%
    \startMPinclusions
        def leftangle =
            pointC ;
            jasper_curve[
                ustart="0"
                ,ustop="pi/2"
                ,fx="cx+2*cos(pi-u)"
                ,fy="cy-1+2*sin(pi-u)"
                ,drawcommands="jasper_thin"
            ] ;
            jasper_curve[
                ustart="0.8"
                ,ustop="1.2"
                ,fx="cx+u*2*cos(pi-(pi/4))"
                ,fy="cy-1+u*2*sin(pi-(pi/4))"
                ,drawcommands="jasper_thin"
            ] ;
        enddef ;
    \stopMPinclusions%
    % build angle NCA
    \dorecurse{4*30}{%
        \startMPpage
            leftangle ;
            jasper_curve[
                ustart="0"
                ,ustop="pi/2"
                ,fx="cx+2*cos(pi-u-(pi/2)*#1/120)"
                ,fy="cy-1+2*sin(pi-u-(pi/2)*#1/120)"
                ,fz="5"
                ,drawcommands="jasper_lightgraypenthin"
            ] ;
            jasper_curve[
                ustart="0.8"
                ,ustop="1.2"
                ,fx="cx+u*2*cos(pi-(pi/4)-(pi/2)*#1/120)"
                ,fy="cy-1+u*2*sin(pi-(pi/4)-(pi/2)*#1/120)"
                ,drawcommands="jasper_lightgraypenthin"
                ,fz="5"
            ] ;
            jasper_render_segments ;
        \stopMPpage%
    }%
    \startMPinclusions
        def rightangle =
            leftangle ;
            jasper_curve[
                ustart="0"
                ,ustop="pi/2"
                ,fx="cx+2*cos(pi-u-(pi/2))"
                ,fy="cy-1+2*sin(pi-u-(pi/2))"
                ,fz="1"
                ,drawcommands="jasper_lightgraypenthin"
            ] ;
            jasper_curve[
                ustart="0.8"
                ,ustop="1.2"
                ,fx="cx+u*2*cos(pi-(pi/4)-(pi/2))"
                ,fy="cy-1+u*2*sin(pi-(pi/4)-(pi/2))"
                ,drawcommands="jasper_lightgraypenthin"
                ,fz="1"
            ] ;
        enddef ;
    \stopMPinclusions%
    % pause after building NCA
    \dorecurse{4*30}{%
        \startMPpage
            rightangle ;
            jasper_render_segments ;
        \stopMPpage%
    }%
    \startbuffer[BufA]
        \startframed[myframed]
            {\bf Theorem 3.2.1} There exists an angle congruent to
            its supplement.\par 
            {\it Proof.} Let \im{AB} be a straight line and \im{M}
            be a point not lying on \im{AB}. Let us subject the 
            plane to a type (iii) isometry: reflection in \im{AB}.
            That means the isometry will carry every point on one 
            side of \im{AB} into a point lying on the other side 
            of this line, and every point of \im{AB} will remain 
            stationary.\par 
            As a result of this transformation, point \im{M} will 
            be carried into some point \im{N} lying on the other 
            side of \im{AB}.\par
            Let us join \im{M} with \im{N} by means of a straight 
            line. This line will intersect \im{AB} at 
            some point, which we shall denote \im{C}.\par 
            Since \im{CA} does not move as a result of the isometry
            performed, \im{\angle MCA} is transformed into 
            \im{\angle NCA}, which means that these two angles 
            are congruent: \im{\angle MCA = \angle NCA}. \color[white]{At the 
            same time, since the two sides (\im{MC} and \im{CN})
            of these angles extend each other into a straight line
            (\im{MN}), these angles are supplementary. Therefore
            \im{\angle MCA} and \im{\angle NCA} are congruent 
            supplementary angles, which means such angles exist,
            {\bf Q.E.D.}}
        \stopframed
    \stopbuffer
    % which means that these two angles 
    % are congruent: \im{\angle MCA = \angle NCA}.\
    \dorecurse{8*30}{%
        \startMPpage
            rightangle ;
            jasper_render_segments ;
        \stopMPpage%
    }%
    \startbuffer[BufA]
        \startframed[myframed]
            {\bf Theorem 3.2.1} There exists an angle congruent to
            its supplement.\par 
            {\it Proof.} Let \im{AB} be a straight line and \im{M}
            be a point not lying on \im{AB}. Let us subject the 
            plane to a type (iii) isometry: reflection in \im{AB}.
            That means the isometry will carry every point on one 
            side of \im{AB} into a point lying on the other side 
            of this line, and every point of \im{AB} will remain 
            stationary.\par 
            As a result of this transformation, point \im{M} will 
            be carried into some point \im{N} lying on the other 
            side of \im{AB}.\par
            Let us join \im{M} with \im{N} by means of a straight 
            line. This line will intersect \im{AB} at 
            some point, which we shall denote \im{C}.\par 
            Since \im{CA} does not move as a result of the isometry
            performed, \im{\angle MCA} is transformed into 
            \im{\angle NCA}, which means that these two angles 
            are congruent: \im{\angle MCA = \angle NCA}. At the 
            same time, since the two sides (\im{MC} and \im{CN})
            of these angles extend each other into a straight line
            (\im{MN}), these angles are supplementary. \color[white]{Therefore
            \im{\angle MCA} and \im{\angle NCA} are congruent 
            supplementary angles, which means such angles exist,
            {\bf Q.E.D.}}
        \stopframed
    \stopbuffer
    % At the 
    % same time, since the two sides (\im{MC} and \im{CN})
    % of these angles extend each other into a straight line
    % (\im{MN}), these angles are supplementary.
    \dorecurse{15*30}{%
        \startMPpage
            rightangle ;
            jasper_render_segments ;
        \stopMPpage%
    }%
    \startbuffer[BufA]
        \startframed[myframed]
            {\bf Theorem 3.2.1} There exists an angle congruent to
            its supplement.\par 
            {\it Proof.} Let \im{AB} be a straight line and \im{M}
            be a point not lying on \im{AB}. Let us subject the 
            plane to a type (iii) isometry: reflection in \im{AB}.
            That means the isometry will carry every point on one 
            side of \im{AB} into a point lying on the other side 
            of this line, and every point of \im{AB} will remain 
            stationary.\par 
            As a result of this transformation, point \im{M} will 
            be carried into some point \im{N} lying on the other 
            side of \im{AB}.\par
            Let us join \im{M} with \im{N} by means of a straight 
            line. This line will intersect \im{AB} at 
            some point, which we shall denote \im{C}.\par 
            Since \im{CA} does not move as a result of the isometry
            performed, \im{\angle MCA} is transformed into 
            \im{\angle NCA}, which means that these two angles 
            are congruent: \im{\angle MCA = \angle NCA}. At the 
            same time, since the two sides (\im{MC} and \im{CN})
            of these angles extend each other into a straight line
            (\im{MN}), these angles are supplementary. Therefore
            \im{\angle MCA} and \im{\angle NCA} are congruent 
            supplementary angles, \color[white]{which means such angles exist,
            {\bf Q.E.D.}}
        \stopframed
    \stopbuffer
    % Therefore
    % \im{\angle MCA} and \im{\angle NCA} are congruent 
    % supplementary angles,
    \dorecurse{10*30}{%
        \startMPpage
            rightangle ;
            jasper_render_segments ;
        \stopMPpage%
    }%
    \startbuffer[BufA]
        \startframed[myframed]
            {\bf Theorem 3.2.1} There exists an angle congruent to
            its supplement.\par 
            {\it Proof.} Let \im{AB} be a straight line and \im{M}
            be a point not lying on \im{AB}. Let us subject the 
            plane to a type (iii) isometry: reflection in \im{AB}.
            That means the isometry will carry every point on one 
            side of \im{AB} into a point lying on the other side 
            of this line, and every point of \im{AB} will remain 
            stationary.\par 
            As a result of this transformation, point \im{M} will 
            be carried into some point \im{N} lying on the other 
            side of \im{AB}.\par
            Let us join \im{M} with \im{N} by means of a straight 
            line. This line will intersect \im{AB} at 
            some point, which we shall denote \im{C}.\par 
            Since \im{CA} does not move as a result of the isometry
            performed, \im{\angle MCA} is transformed into 
            \im{\angle NCA}, which means that these two angles 
            are congruent: \im{\angle MCA = \angle NCA}. At the 
            same time, since the two sides (\im{MC} and \im{CN})
            of these angles extend each other into a straight line
            (\im{MN}), these angles are supplementary. Therefore
            \im{\angle MCA} and \im{\angle NCA} are congruent 
            supplementary angles, which means such angles exist,
            \color[white]{{\bf Q.E.D.}}
        \stopframed
    \stopbuffer
    % which means such angles exist,
    \dorecurse{5*30}{%
        \startMPpage
            rightangle ;
            jasper_render_segments ;
        \stopMPpage%
    }%
    \startbuffer[BufA]
        \startframed[myframed]
            {\bf Theorem 3.2.1} There exists an angle congruent to
            its supplement.\par 
            {\it Proof.} Let \im{AB} be a straight line and \im{M}
            be a point not lying on \im{AB}. Let us subject the 
            plane to a type (iii) isometry: reflection in \im{AB}.
            That means the isometry will carry every point on one 
            side of \im{AB} into a point lying on the other side 
            of this line, and every point of \im{AB} will remain 
            stationary.\par 
            As a result of this transformation, point \im{M} will 
            be carried into some point \im{N} lying on the other 
            side of \im{AB}.\par
            Let us join \im{M} with \im{N} by means of a straight 
            line. This line will intersect \im{AB} at 
            some point, which we shall denote \im{C}.\par 
            Since \im{CA} does not move as a result of the isometry
            performed, \im{\angle MCA} is transformed into 
            \im{\angle NCA}, which means that these two angles 
            are congruent: \im{\angle MCA = \angle NCA}. At the 
            same time, since the two sides (\im{MC} and \im{CN})
            of these angles extend each other into a straight line
            (\im{MN}), these angles are supplementary. Therefore
            \im{\angle MCA} and \im{\angle NCA} are congruent 
            supplementary angles, which means such angles exist,
            {\bf Q.E.D.}
        \stopframed
    \stopbuffer
    % {\bf Q.E.D.}
    \dorecurse{10*30}{%
        \startMPpage
            rightangle ;
            jasper_render_segments ;
        \stopMPpage%
    }%
\stoptext
